\begin{longtable}[c]{|c|m{0.16\textwidth}|m{0.43\textwidth}|m{0.21\textwidth}|}
    \caption{Épicas de Usuario del proyecto.}
    \label{tab:epicasDeUsuario} \\
    \hline
    \rowcolor{gray!50}
    \multicolumn{1}{|c|}{\textbf{ID}} & \multicolumn{1}{c|}{\textbf{Épica}} & \multicolumn{1}{c|}{\textbf{Descripción}} & \multicolumn{1}{c|}{\textbf{HUs}} \\
    \hline
    \endfirsthead
    
    \hline
    \rowcolor{gray!50}
    \multicolumn{1}{|c|}{\textbf{ID}} & \multicolumn{1}{c|}{\textbf{Épica}} & \multicolumn{1}{c|}{\textbf{Descripción}} & \multicolumn{1}{c|}{\textbf{HUs}} \\
    \hline
    \endhead
    
    \multicolumn{4}{r}{\textit{Continúa en la siguiente página...}} \\
    \endfoot
    
    \endlastfoot
    
    AUM & Autenticación y gestión de usuarios & Facilita a todo desarrollador registrarse e iniciar sesión. Además brinda a los administradores, el control sobre los usuarios registrados para otorgar o denegar permisos dentro del sistema. & AUM-01, AUM-02, AUM-03 y AUM-04 \\
    \hline
    CON & Gestión de conexiones & Permite al usuario configurar las conexiones a los DBMS SQL Server y Neo4j para la interacción con las bases de datos correspondientes. & CON-01 y CON-02 \\
    \hline
    QTE & Traducción y ejecución de consultas & Se encarga del análisis, traducción de sentencias SQL a lenguaje Cypher y su posterior ejecución en una base de datos orientada a grafos. También incluye el historial de consultas y la visualización de resultados obtenidos en distintos formatos. & QTE-01, QTE-02, QTE-03 y QTE-04 \\
    \hline
    OBS & Observabilidad y seguimiento & Permite a los administradores la realización de auditoría y monitoreo de rendimiento para identificar áreas de mejora. & OBS-01 y OBS-02 \\
    \hline
    
\end{longtable}